\subsection{Adapting the Tool to Bitcoin}
\label{sec:adaptation}

The pipeline above was written for the Python corpus and ran on the Bitcoin corpus after four changes, none of which touches the classifiers. First, the table and column names that identify a proposal are derived from the proposal identifier (\texttt{accrejbips}, \texttt{bipdetails}, \texttt{bipeditors}, column \texttt{BIP}), so the same code reads either database. Second, the outcome resolver understands the BIP status vocabulary, including the 2025 revision (Section~\ref{sec:corpus}). Third, the role mapper knows the \emph{lead maintainer} role, which takes the BDFL's weight in the score, and reads Bitcoin's roster tables; the message-level role labels are computed by a post-reading script from those rosters exactly as for Python. Fourth, the cue lists were extended with Bitcoin's own governance vocabulary---its decision-makers (\emph{lead maintainer, core maintainers, BIP editors} and the handles of the best-known maintainers) for \emph{authority}, the companies and funders active in the project for \emph{corporate interest}, chain splits and client switching for \emph{exit threat}, list moderation and censorship complaints for \emph{procedural control}, and the flame vocabulary of the block-size war (\emph{shill, scam, sockpuppet, brigading, bad faith}) for \emph{incivility} (Appendix~\ref{app:rules}). Release announcements, whose release notes enumerate dozens of BIPs and are not discussion, are skipped by a subject filter, as commit notifications were skipped for Python; without it the lead maintainer's release notes alone would have contributed more than 5,000 candidate sentences. Both user interfaces read the Bitcoin database unchanged.

